% The `physics` package (CTAN, Sergio C. de la Barrera), as plain macros.
% Names and output follow the package manual; see
% https://ctan.org/pkg/physics.
%
% Every renderer here reads the same `\newcommand` subset (KaTeX in the page,
% RaTeX in Emacs, the generated .sty in LaTeX export), so only commands with a
% fixed argument count are ported. Left out on purpose:
%   - delimiter-scanning forms: \qty( ), \mqty( ), \eval( |, \sin( ), \dd( )
%     — use \pqty{} \bqty{} \Bqty{} \vqty{} and \pmqty{} \bmqty{} \vmqty{};
%   - starred and optional-power forms (\dv[n], \abs*, \qq*);
%   - redefinitions of built-ins: \div \Re \Im \sin \exp \log \det \Pr;
%   - names other shared files already define with their own meaning:
%     \abs \norm \bra \ket \braket \ketbra \expval \ip \tr \Tr \rank \vec.
%   - computed matrices: \imat \xmat \zmat \pmat \dmat \admat.

% 2.1 Automatic bracing
\newcommand{\pqty}[1]{\left(#1\right)}
\newcommand{\bqty}[1]{\left[#1\right]}
\newcommand{\Bqty}[1]{\left\{#1\right\}}
\newcommand{\vqty}[1]{\left|#1\right|}
\newcommand{\absolutevalue}[1]{\left|#1\right|}
\newcommand{\evaluated}[1]{\left.#1\right|}
\newcommand{\eval}[1]{\left.#1\right|}
\newcommand{\order}[1]{\mathcal{O}\left(#1\right)}
\newcommand{\commutator}[2]{\left[#1,#2\right]}
\newcommand{\comm}[2]{\left[#1,#2\right]}
\newcommand{\anticommutator}[2]{\left\{#1,#2\right\}}
\newcommand{\acomm}[2]{\left\{#1,#2\right\}}
\newcommand{\poissonbracket}[2]{\left\{#1,#2\right\}}
\newcommand{\pb}[2]{\left\{#1,#2\right\}}

% 2.2 Vector notation
\newcommand{\vectorbold}[1]{\mathbf{#1}}
\newcommand{\vb}[1]{\mathbf{#1}}
\newcommand{\vectorarrow}[1]{\overrightarrow{\mathbf{#1}}}
\newcommand{\va}[1]{\overrightarrow{\mathbf{#1}}}
\newcommand{\vectorunit}[1]{\hat{\mathbf{#1}}}
\newcommand{\vu}[1]{\hat{\mathbf{#1}}}
\newcommand{\dotproduct}{\cdot}
\newcommand{\vdot}{\cdot}
\newcommand{\crossproduct}{\times}
\newcommand{\cross}{\times}
\newcommand{\cp}{\times}
\newcommand{\gradient}{\nabla}
\newcommand{\grad}{\nabla}
\newcommand{\divergence}{\nabla\cdot}
\newcommand{\curl}{\nabla\times}
\newcommand{\laplacian}{\nabla^2}

% 2.3 Operators
\newcommand{\trace}{\operatorname{tr}}
\newcommand{\Trace}{\operatorname{Tr}}
\newcommand{\erf}{\operatorname{erf}}
\newcommand{\Res}{\operatorname{Res}}
\newcommand{\principalvalue}{\mathcal{P}}
\newcommand{\pv}{\mathcal{P}}
\newcommand{\PV}{\mathrm{P.V.}}

% 2.4 Quick quad text
\newcommand{\qqtext}[1]{\quad\text{#1}\quad}
\newcommand{\qq}[1]{\quad\text{#1}\quad}
\newcommand{\qcomma}{,\quad}
\newcommand{\qc}{,\quad}
\newcommand{\qcc}{\quad\text{c.c.}\quad}
\newcommand{\qif}{\quad\text{if}\quad}
\newcommand{\qthen}{\quad\text{then}\quad}
\newcommand{\qelse}{\quad\text{else}\quad}
\newcommand{\qotherwise}{\quad\text{otherwise}\quad}
\newcommand{\qunless}{\quad\text{unless}\quad}
\newcommand{\qgiven}{\quad\text{given}\quad}
\newcommand{\qusing}{\quad\text{using}\quad}
\newcommand{\qassume}{\quad\text{assume}\quad}
\newcommand{\qsince}{\quad\text{since}\quad}
\newcommand{\qlet}{\quad\text{let}\quad}
\newcommand{\qfor}{\quad\text{for}\quad}
\newcommand{\qall}{\quad\text{all}\quad}
\newcommand{\qeven}{\quad\text{even}\quad}
\newcommand{\qodd}{\quad\text{odd}\quad}
\newcommand{\qinteger}{\quad\text{integer}\quad}
\newcommand{\qand}{\quad\text{and}\quad}
\newcommand{\qor}{\quad\text{or}\quad}
\newcommand{\qas}{\quad\text{as}\quad}
\newcommand{\qin}{\quad\text{in}\quad}

% 2.5 Derivatives (two-argument forms)
\newcommand{\differential}[1]{\mathop{}\!\mathrm{d}#1}
\newcommand{\dd}[1]{\mathop{}\!\mathrm{d}#1}
\newcommand{\derivative}[2]{\frac{\mathrm{d}#1}{\mathrm{d}#2}}
\newcommand{\dv}[2]{\frac{\mathrm{d}#1}{\mathrm{d}#2}}
\newcommand{\partialderivative}[2]{\frac{\partial#1}{\partial#2}}
\newcommand{\pderivative}[2]{\frac{\partial#1}{\partial#2}}
\newcommand{\pdv}[2]{\frac{\partial#1}{\partial#2}}
\newcommand{\variation}[1]{\mathop{}\!\delta#1}
\newcommand{\var}[1]{\mathop{}\!\delta#1}
\newcommand{\functionalderivative}[2]{\frac{\delta#1}{\delta#2}}
\newcommand{\fdv}[2]{\frac{\delta#1}{\delta#2}}

% 2.6 Dirac notation (beyond quantum.tex)
\newcommand{\innerproduct}[2]{\left\langle#1\middle|#2\right\rangle}
\newcommand{\outerproduct}[2]{\left|#1\right\rangle\!\left\langle#2\right|}
\newcommand{\dyad}[2]{\left|#1\right\rangle\!\left\langle#2\right|}
\newcommand{\op}[2]{\left|#1\right\rangle\!\left\langle#2\right|}
\newcommand{\expectationvalue}[1]{\left\langle#1\right\rangle}
\newcommand{\ev}[1]{\left\langle#1\right\rangle}
\newcommand{\matrixelement}[3]{\left\langle#1\middle|#2\middle|#3\right\rangle}
\newcommand{\matrixel}[3]{\left\langle#1\middle|#2\middle|#3\right\rangle}
\newcommand{\mel}[3]{\left\langle#1\middle|#2\middle|#3\right\rangle}

% 2.7 Matrix macros
\newcommand{\matrixquantity}[1]{\begin{matrix}#1\end{matrix}}
\newcommand{\mqty}[1]{\begin{matrix}#1\end{matrix}}
\newcommand{\pmqty}[1]{\begin{pmatrix}#1\end{pmatrix}}
\newcommand{\Pmqty}[1]{\left(\begin{matrix}#1\end{matrix}\right)}
\newcommand{\bmqty}[1]{\begin{bmatrix}#1\end{bmatrix}}
\newcommand{\vmqty}[1]{\begin{vmatrix}#1\end{vmatrix}}
\newcommand{\smallmatrixquantity}[1]{\begin{smallmatrix}#1\end{smallmatrix}}
\newcommand{\smqty}[1]{\begin{smallmatrix}#1\end{smallmatrix}}
\newcommand{\spmqty}[1]{\left(\begin{smallmatrix}#1\end{smallmatrix}\right)}
\newcommand{\sPmqty}[1]{\left(\begin{smallmatrix}#1\end{smallmatrix}\right)}
\newcommand{\sbmqty}[1]{\left[\begin{smallmatrix}#1\end{smallmatrix}\right]}
\newcommand{\svmqty}[1]{\left|\begin{smallmatrix}#1\end{smallmatrix}\right|}
\newcommand{\matrixdeterminant}[1]{\begin{vmatrix}#1\end{vmatrix}}
\newcommand{\mdet}[1]{\begin{vmatrix}#1\end{vmatrix}}
\newcommand{\smdet}[1]{\left|\begin{smallmatrix}#1\end{smallmatrix}\right|}
